\section{Acquisition reward, duplicate consumption and forecast validation}
\label{sec:acquisition-supplement}
These diagnostics distinguish hypothetical visible area, actual gross-new object area, and net area after native TSDF fusion. Gross area is evaluator source surface farther than 5\,cm from the before mesh and within 5\,cm of the after mesh. Lost area has the reverse transition, and net equals gross minus lost. Whole-scene metrics use the complete reference; object metrics use the two modeled object subsets. Alternative-template quadrature weights are expected areas, not a union of independent physical surfaces. All 66 existing branches and eight mechanism branches are retained; repeated predictions do not create additional independent acquisitions.

\subsection{Measured-mesh duplication and complete-policy ablation}
Across the six retained prefixes, the G/S weighted reward fractions are respectively 42.78/39.75\% unsupported by object source geometry, 6.05/6.41\% object-supported but already reconstructed, 26.41/31.08\% supported and newly realized, and 24.75/22.77\% supported but unrealized after the branch. Fractions first average candidate area within each prefix and then average prefixes equally before normalization. Including points unsupported by the object reference, 7.67/7.48\% of total expected reward is already within 5\,cm of the native mesh. Another 4.67/5.58\% of total reward lies near non-object scene geometry despite lacking object support. Reference proximity does not by itself establish visible external surface.

The additive native mask sets the existing reward to zero where current native distance is at most 5\,cm. It leaves template probabilities, original history masks, candidates, common scores and execution unchanged. This is a declared reward revision, shared by G/S. It removes the measured duplicate proxy, but can also consume points near erroneous reconstructed geometry. Table~\ref{tab:gain-replay} retains the slight G rank degradation and all unchanged winners. Four complete 256-action runs then remain exactly identical in trajectory and native mesh to their respective $T=16$ controls (Table~\ref{tab:native-consumption-policy}). Full coverage, zero simulated collisions and return are retained; no improvement in exploration or reconstruction is observed. The adverse left semantic comparison remains unchanged.
\begin{table}[!htbp]
\centering\small
\caption{All 66 reused candidate branches at six development prefixes. Errors compare predicted object area with actual object gross or net newly reconstructed area at 5\,cm; candidates are averaged within each prefix, then prefixes equally. Smaller MAE and larger rank correlation are preferable. All twelve prefix--policy winners remain unchanged.}
\label{tab:gain-replay}
\begin{tabular}{llrrrr}
\toprule
Prior & Consumption & Gross MAE & Net MAE & Gross $\rho$ & Net $\rho$ \\
\midrule
G & History & 0.2413 & 0.3407 & 0.9272 & 0.8648 \\
G & History + native & 0.2304 & 0.3141 & 0.9189 & 0.8596 \\
S & History & 0.2566 & 0.3489 & 0.9161 & 0.8556 \\
S & History + native & 0.2447 & 0.3265 & 0.9184 & 0.8580 \\
\bottomrule
\end{tabular}
\end{table}

\begin{table}[!htbp]
\centering\small
\caption{Four complete 256-action native-consumption runs on the same two development layouts. Full denotes the first action at 100\% correct-free coverage. F1 and object recalls are percentages at 5\,cm. The final column tests exact trajectory and native-mesh identity with the corresponding frozen $T=16$ control. Three externally interrupted partial attempts are retained separately and excluded from complete-run rows.}
\label{tab:native-consumption-policy}
\begin{tabular}{llrrrrrl}
\toprule
Layout & Prior & Full & Path (m) & F1 & Shelf & Cabinet & Same as control \\
\midrule
Left & G & 88 & 42.0 & 92.758 & 82.170 & 88.709 & Yes \\
Left & S & 97 & 43.2 & 92.381 & 83.464 & 85.874 & Yes \\
Right & G & 100 & 43.2 & 92.596 & 81.088 & 79.214 & Yes \\
Right & S & 100 & 43.2 & 92.596 & 81.088 & 79.214 & Yes \\
\bottomrule
\end{tabular}
\end{table}


The first batch was terminated externally with three incomplete attempts at 240, 238 and 202 actions and the fourth not started. These partial outputs and interruption receipt remain archived. A detached rerun used the same frozen source/configuration to obtain the four complete rows; interrupted fragments are neither complete trials nor independent repeats. Original wall durations of the complete runs are 423.65--500.94\,s under up to three concurrent processes, with numerical-library thread limits of one. They are shared-host computation measurements, not controlled real-time or physical travel estimates.

\subsection{Known-geometry consumption probes and failed alternatives}
The first twelve objects contain six closed and six open structures with independently sampled dimensions. Each has three fixed routes after three acquired prefix frames: revisit, side and around. Geometry is known to the hypothetical model; independent source Monte Carlo samples provide the actual gross/lost/net area. The additive native mask worsens overall MAE, improving three route errors, worsening ten and leaving 23 tied. Two subsequent development variants replace the old 10\,cm measured-point consumption with native 5\,cm consumption, first on the original coarse quadrature and then on a finer basis. Both also worsen overall error (Table~\ref{tab:gain-fixture-ablation}). In particular, a point absent from the present native mesh may still fail to be reconstructed by a repeated view. The additive mask is not promoted to the default reward.
\begin{table}[!htbp]
\centering\small
\caption{Known-geometry acquisition probes on twelve objects (six closed, six open), each with three fixed routes. MAE and signed error are in m$^2$ against gross newly reconstructed source area. Append native was specified before these observations; replacement variants reuse the same objects as development diagnostics. These are sensor/fusion probes, not inferred-geometry or complete-exploration tests.}
\label{tab:gain-fixture-ablation}
\begin{tabular}{lrrrrr}
\toprule
Reward accounting & All MAE & Revisit & Side & Around & Signed error \\
\midrule
Original & 0.1563 & 0.0379 & 0.1271 & 0.3040 & -0.1311 \\
Append native & 0.1669 & 0.0253 & 0.1432 & 0.3321 & -0.1500 \\
Replace history & 0.1865 & 0.2510 & 0.1823 & 0.1262 & +0.1542 \\
Replace + finer quadrature & 0.2313 & 0.3171 & 0.2088 & 0.1681 & +0.2073 \\
\bottomrule
\end{tabular}
\end{table}


\subsection{Forward acquisition operator and independent object split}
The forward estimator uses fresh prediction-only TSDF volumes, 6\,cm voxels, 12\,cm truncation and the same acquired depth convention. All past frames are retained; cached meshes are not treated as substitutes for TSDF integration weights. Hypothetical future depth never updates the real measured map. Its own surface basis uses 0.06\,m quadrature, outward primitive-boundary checks, coincident-face area splitting and twelve fixed orbital first-hit views. This is a discrete modeled surface basis, not an exact analytic union. The source evaluator independently uses 30,000 Monte Carlo samples per object and does not supply points to prediction.

The estimator and twelve new test objects were frozen before outcomes. The earlier twelve objects are development data, not another independent test. Each split has six closed and six open objects and 36 dependent routes. All three predictors in Table~\ref{tab:acquisition-fixtures} share this fine surface basis: old-history visibility, native-consumption replacement visibility, and forward acquisition. Native replacement is different from the additive mask above. All route predictions are completed before source-reference and actual-future evaluation for that object. With known exact geometry, independently rendered hypothetical and actual future depth, as well as resulting native mesh vertices/triangles, match exactly. This is an acquisition/fusion consistency probe rather than evidence of successful unseen-shape prediction. Frozen toy render helpers have slightly different range-boundary conventions; exact depth-array assertions verify no difference in these evaluated views.

Forward MAE is 0.0191\,m$^2$ on development objects and 0.0143\,m$^2$ on new test objects. Against the matched visibility control, development has 30 improved, one worsened and five tied route errors; the test has 33 improved, one worsened and two tied errors. Revisit gross gain is zero in the actual reconstruction and in the forward estimate, while test visibility predicts 0.1006\,m$^2$ on average. The remaining discrepancies arise from the independent quadratures even though the known-geometry forecast and actual native meshes match. This favorable result is not used as a substitute for fitted-template or complete-policy evidence.
\begin{table}[!htbp]
\centering\small
\caption{Matched 0.06\,m surface-basis probes: twelve objects per split, six closed and six open, with three dependent routes each. Units are m$^2$. Native replacement substitutes a 5\,cm native-mesh consumption rule for the old 10\,cm measured-point rule; it is not the additive mask used in the complete-policy ablation. All predictors use the same known geometry and source evaluation.}
\label{tab:acquisition-fixtures}
\begin{tabular}{llrrrrr}
\toprule
Split & Estimator & All MAE & Revisit & Side & Around & Bias \\
\midrule
Development & Visibility / old history & 0.1578 & 0.0607 & 0.1298 & 0.2828 & -0.1173 \\
Development & Visibility / native replacement & 0.2268 & 0.3135 & 0.2021 & 0.1649 & +0.2009 \\
Development & Acquisition forecast & 0.0191 & 0.0000 & 0.0244 & 0.0330 & +0.0018 \\
New test & Visibility / old history & 0.1532 & 0.1006 & 0.1321 & 0.2270 & -0.0664 \\
New test & Visibility / native replacement & 0.2228 & 0.3192 & 0.2019 & 0.1474 & +0.2011 \\
New test & Acquisition forecast & 0.0143 & 0.0000 & 0.0275 & 0.0153 & +0.0015 \\
\bottomrule
\end{tabular}
\end{table}


\subsection{Fitted geometry at all original action-54 candidates}
The fitted probe replays exactly 55 acquired RGB-D frames and retains the original 21 history keyframes at the left-layout prefix after 54 actions. It exports both structures for each of two observed instances at the class-independent conditional mean of depth and placement. This is an explicit point approximation to the original conditional geometry distribution. For each of the eleven original outbound routes, future images contain the current acquired mesh and exactly one alternative, giving 44 hypothesis--route forecasts. G and S share these forecasts and differ only by the original $T=16$ structure probabilities. Strict-free contradictions can remove hypothetical first intersections only with the original cross-view measured-support veto; current native intersections are never removed. The fitted renderer preserves the actual noiseless chain's inclusive 0.15--8\,m axial range. Unmodeled unseen occluders remain absent.

All 44 predictions are completed before independent binding to the eleven previously acquired actual branch endpoints. No future actual depth or evaluator source geometry enters prediction. Table~\ref{tab:acquisition-fitted} preserves every comparison. Relative to mean/fine old-history visibility, forward gross MAE rises from 0.3107 to 0.3172\,m$^2$ for G and from 0.3961 to 0.3973\,m$^2$ for S, despite improved object-gross rank correlation. The original full-distribution predictor has smaller error and stronger ranking than the forward approximation. Forecast net-versus-actual-net MAE is 0.3148/0.4894\,m$^2$ for G/S; net is not substituted into the original gross score.

All three mean/fine predictors select candidate 8 for both priors. Its actual object/global net additions are 1.8049/16.3250\,m$^2$ over 51 actions and 8.8\,m. The original distribution's S selects candidate 0, with 1.8482/17.3831\,m$^2$ over 55 actions and 9.4\,m; free-coverage increments are identical. The change from candidate 0 to 8 therefore occurs already under the mean/basis approximation and cannot be attributed to forward acquisition. This single development prefix does not establish generalization or a complete-policy improvement. That conditional-mean/fine-basis variant remains a historical prediction probe. The new full-posterior, original-quadrature complete-policy comparison is reported in the main text and the following four-layout section.
\begin{table}[!htbp]
\centering\small\setlength{\tabcolsep}{3pt}
\caption{All eleven reused actual branch outcomes at left-prefix action 54. Gross prediction error and its mismatch to actual net area are both retained; the latter does not redefine the score. Selected object/global net area is measured independently. All mean/fine variants select candidate 8; selection already changes before replacing visibility with acquisition forecasting.}
\label{tab:acquisition-fitted}
\begin{tabular}{llrrrrrrr}
\toprule
Estimator & Prior & Gross MAE & Net mismatch & $\rho$ & Choice & Obj. net & Global net & Actions \\
\midrule
Distribution & G & 0.2061 & 0.1943 & 0.927 & 8 & 1.8049 & 16.325 & 51 \\
Distribution & S & 0.1961 & 0.1715 & 0.891 & 0 & 1.8482 & 17.383 & 55 \\
Mean / old history & G & 0.3107 & 0.2619 & 0.691 & 8 & 1.8049 & 16.325 & 51 \\
Mean / old history & S & 0.3961 & 0.3270 & 0.791 & 8 & 1.8049 & 16.325 & 51 \\
Mean / native replacement & G & 0.3070 & 0.2614 & 0.700 & 8 & 1.8049 & 16.325 & 51 \\
Mean / native replacement & S & 0.3563 & 0.2878 & 0.764 & 8 & 1.8049 & 16.325 & 51 \\
Mean / acquisition & G & 0.3172 & 0.2663 & 0.836 & 8 & 1.8049 & 16.325 & 51 \\
Mean / acquisition & S & 0.3973 & 0.3279 & 0.855 & 8 & 1.8049 & 16.325 & 51 \\
\bottomrule
\end{tabular}
\end{table}

